% ID: SEL-D0044-OA-E02
% ESTADO: validada. S045, cotejada con examen y criterios.
\textbf{a)} $f(x)=(x-1)^3+1$ es estrictamente creciente
en todo $\mathbb R$, porque lo es la función cúbica básica.
También se observa que $f'(x)=3(x-1)^2$ sólo se anula en $1$,
sin cambiar de signo. No hay extremos relativos.

\textbf{b)} $f''(x)=6(x-1)$. La gráfica es cóncava hacia abajo
para $x<1$ y hacia arriba para $x>1$. En $x=1$ cambia la
curvatura y $f(1)=1$, de modo que hay una inflexión con tangente
horizontal en $\boxed{(1,1)}$.

\textbf{c)} La gráfica es la de $x^3$ trasladada una unidad
a la derecha y una hacia arriba. Pasa por $(0,0)$, $(1,1)$
y $(2,2)$.
\begin{center}
\setlength{\unitlength}{7mm}
\begin{picture}(7,9)(-3,-4)
\qbezier(-2.5,0)(0.5,0)(3.5,0)\put(3.6,0){$x$}
\qbezier(0,-3.5)(0,0.5)(0,4.5)\put(0,4.6){$y$}
\qbezier(-0.5,-2.375)(-0.25,-0.71875)(0,0)
\qbezier(0,0)(0.25,0.71875)(0.5,0.875)
\qbezier(0.5,0.875)(0.75,1.03125)(1,1)
\qbezier(1,1)(1.25,0.96875)(1.5,1.125)
\qbezier(1.5,1.125)(1.75,1.28125)(2,2)
\qbezier(2,2)(2.25,2.71875)(2.5,4.375)
\put(1,0.8){\makebox(0,0)[t]{$(1,1)$}}
\end{picture}
\end{center}
